\chapter{LogLikelihood}
\label{ch:loglikelihood}

\newthought{\yamlkey{LogLikelihood}} is how a point gets a score. After the workflow turns a
parameter point into observables, \Jarvis\ evaluates these expressions to decide how good the
point is---and samplers use that score to decide what to do next.

\section{The minimal form}
\label{sec:logl-minimal}

\yamlkey{LogLikelihood} lives inside \yamlkey{Sampling} and is a list of named terms. The
simplest case is a single term:

\begin{yamlwin}
Sampling:
  LogLikelihood:
    - name: LogL
      expression: "-0.5 * chi2"
\end{yamlwin}

\begin{description}[style=nextline, leftmargin=1.2em]
  \item[\yamlkey{name}] A stable name for the term. \Jarvis\ writes its value into the sample
        record under this name, so you can inspect it afterwards.
  \item[\yamlkey{expression}] A symbolic expression (Chapter~\ref{ch:symbols}) evaluated once the
        scan variables and workflow outputs are available.
\end{description}

\begin{jnote}
Every term you list is computed and saved---even intermediate ones---so they are all available
for later analysis, not just the final score.
\end{jnote}

\section{Several terms at once}
\label{sec:logl-multi}

Real fits combine several constraints. List each as its own term so you can inspect them
separately:

\begin{yamlwin}
Sampling:
  LogLikelihood:
    - {name: "LogL_Z", expression: "LogGauss(z, 100, 10)"}
    - {name: "LogL_H", expression: "LogGauss(h, 125, 2)"}
\end{yamlwin}

\noindent Here \yamlkey{z} and \yamlkey{h} are observables produced by the workflow, and
\expr{LogGauss(value, mean, error)} is a built-in Gaussian log-likelihood
(Chapter~\ref{ch:symbols}).

\section{How the total score is formed}
\label{sec:logl-total}

The total log-likelihood is stored in the standard field \yamlkey{LogL}. There are two ways it is
set:

\paragraph{Automatic sum (default)} If you do \emph{not} define a term named \yamlkey{LogL},
\Jarvis\ adds up all your terms and stores the sum as \yamlkey{LogL}. The two-term example above
therefore produces \yamlkey{LogL\_Z}, \yamlkey{LogL\_H}, and
$\yamlkey{LogL} = \yamlkey{LogL\_Z} + \yamlkey{LogL\_H}$.

\paragraph{Explicit total} If you \emph{do} define a term named \yamlkey{LogL}, that expression
is used as the total. The other terms are still computed and saved, and your \yamlkey{LogL}
expression may reference them by name:

\begin{yamlwin}
Sampling:
  LogLikelihood:
    - {name: "LogL_Z", expression: "LogGauss(z, 100, 10)"}
    - {name: "LogL_H", expression: "LogGauss(h, 125, 2)"}
    - {name: "LogL",   expression: "LogL_Z + 0.5 * LogL_H"}
\end{yamlwin}

\noindent Use the explicit form when terms should be weighted differently, or combined in any way
other than a plain sum.

\begin{jtip}
Keep a single, clearly-named total. Downstream tools (and \JPlot) expect one \yamlkey{LogL}
column as the overall score; the named sub-terms are there for diagnostics.
\end{jtip}

\section{What an expression may use}
\label{sec:logl-inputs}

A \yamlkey{LogLikelihood} expression can reference:

\begin{itemize}
  \item scan variables from \yamlkey{Sampling.Variables} (by their \yamlkey{name});
  \item observables that the workflow (calculators or operas) produced;
  \item built-in constants and functions, and your own \yamlkey{Utils} functions
        (Chapter~\ref{ch:symbols} and Chapter~\ref{ch:utils}).
\end{itemize}

\begin{jwarn}
Every name in an expression must already exist when the term is evaluated. If a likelihood uses
an observable, make sure the module that produces it runs earlier in the workflow---otherwise the
name is undefined.
\end{jwarn}

\section{Summary}
\label{sec:logl-summary}

\yamlkey{LogLikelihood} is a list of \yamlkey{name}/\yamlkey{expression} terms. Omit a
\yamlkey{LogL} term to have \Jarvis\ sum everything automatically, or define \yamlkey{LogL}
explicitly to control the combination. Expressions read scan variables, workflow observables, and
the functions in Chapter~\ref{ch:symbols}.
